\documentclass[11pt,a4paper]{article}
\usepackage[margin=2.2cm]{geometry}
\usepackage[utf8]{inputenc}
\usepackage[T1]{fontenc}
\usepackage{amsmath,amssymb,amsfonts,bm}
\usepackage{braket}
\usepackage{enumitem}
\usepackage{xcolor}
\usepackage{tikz}
\usetikzlibrary{quantikz2}

\newcommand{\pts}[1]{\textbf{(#1)}}
\newcommand{\tr}{\operatorname{tr}}
\newcommand{\id}{\mathbb{I}}
\newcommand{\sol}{\color{blue!60!black}}

\setlist[enumerate,1]{label=\alph*)}
\setlength{\parindent}{0pt}
\setlength{\parskip}{4pt}

\title{Quantum Technologies -- First partial exam\\[4pt] \large Practice exam C}
\date{}

\begin{document}
\maketitle
\vspace{-1.5cm}
\begin{center}
Duration: 2 hours. Total: 10 points.
\end{center}

%=====================================================================
\section{Open question (1pt)}

The quantum Fourier transform (QFT) can be implemented on $n$ qubits with $\mathcal O(n^2)$ gates,
while the classical fast Fourier transform of a vector of $N=2^n$ components needs
$\mathcal O(N\log N)$ operations. Does this mean that a quantum computer can compute the Fourier
transform of classical data exponentially faster? Explain why or why not. Then explain how the QFT
is used in ONE of the algorithms seen in class, and which information it allows us to extract.
(Maximum 200 words, $\sim$1 page)

%=====================================================================
\section{Short questions (3pts)}

\begin{enumerate}
\item Consider the following circuit acting on four qubits. Can it be used as the oracle of the
Grover algorithm? If so, indicate (i) which are the marked items, (ii) the number of Grover
iterations that maximizes the probability of finding a marked item, together with that probability,
and (iii) the probability of success if, by mistake, one more iteration is applied \pts{1 p}.
\begin{center}
\begin{quantikz}[row sep=0.25cm, column sep=0.35cm]
\lstick{$\ket{q_0}$} & & \ctrl{1} & & \\
\lstick{$\ket{q_1}$} & \gate{X} & \ctrl{2} & \gate{X} & \\
\lstick{$\ket{q_2}$} & & & & \\
\lstick{$\ket{q_3}$} & & \gate{Z} & &
\end{quantikz}
\end{center}

\item We use quantum phase estimation with $n=3$ qubits in the first register.
\begin{enumerate}[label=(\roman*)]
\item For $U=TS=\begin{pmatrix}1&0\\0&e^{3i\pi/4}\end{pmatrix}$ and for $U^\dagger$, with the eigenstate
$\ket u=\ket1$, what is the measured string $\vec x$, and with which probability?
\item For $V=\begin{pmatrix}1&0\\0&e^{2\pi i/3}\end{pmatrix}$ and $\ket u=\ket1$, which is the most
likely outcome? Give a lower bound for its probability.
\item How many qubits would be needed in the first register to obtain the phase of $V$ with a
precision of $2^{-3}$ and a probability of success of at least $90\%$?
\end{enumerate}
\pts{1 p}

\item Solving the order-finding problem with a quantum computer for $x=2$ and $M=11$, with $n=8$
qubits in the first register, we measure $\vec y_0=\{01001101\}$. Use the continued fraction expansion
to find the order of $x$ modulo $M$, and check the result. In another run we measure
$\vec y_0^{\,\prime}=\{10000000\}$. Which order do the continued fractions suggest in this case? Is
it correct? How could we still obtain the order from this measurement? \pts{1 p}
\end{enumerate}

%=====================================================================
\section*{Problems}

\section{Amplitude damping (3pts)}

The spontaneous decay of a two-level atom, with $\ket0$ the ground state and $\ket1$ the excited
state, is described by the measurement operators
\[
M_0=\begin{pmatrix}1&0\\0&\sqrt{1-\gamma}\end{pmatrix},\qquad
M_1=\begin{pmatrix}0&\sqrt{\gamma}\\0&0\end{pmatrix},\qquad 0\le\gamma\le1 ,
\]
where outcome $1$ corresponds to the detection of an emitted photon.
\begin{enumerate}
\item Show that $\{M_0,M_1\}$ define a generalized measurement and find the associated POVM
elements $E_0,E_1$. Are they projectors? \pts{0.5 p}
\item For a general state $\rho$, calculate the probability of detecting a photon and the state of
the atom after the measurement for each outcome. Interpret the result \pts{0.5 p}.
\item The outcome is not recorded. Write the output state $\rho'$ and show that, if
$\rho=\frac12(\id+\bm r\cdot\bm\sigma)$, the Bloch vector of the output is
$\bm r'=\left(\sqrt{1-\gamma}\,r_x,\ \sqrt{1-\gamma}\,r_y,\ \gamma+(1-\gamma)r_z\right)$ \pts{0.5 p}.
\item Which state is left invariant by the channel? What happens to an arbitrary state if the
channel is applied many times? What is the output for the maximally mixed state $\id/2$? Compare
with the Pauli channels: can the amplitude damping channel increase the purity of a state? \pts{0.5 p}
\item The input state is $\ket+$. Calculate the purity of the output and its eigenvalues. Evaluate
the von Neumann entropy for $\gamma=1/2$ \pts{0.5 p}.
\item For the input $\ket+$, calculate the fidelity $F(\ket+\!\bra+,\rho')$ and the trace distance
$d(\ket+\!\bra+,\rho')$ between input and output, and evaluate them for $\gamma=1/2$. Compare with
the inputs $\ket0$ and $\ket1$ \pts{0.5 p}.
\end{enumerate}

\section{A rotated Bell pair (3pts)}

Consider the following two-qubit circuit, with
$R_y(\theta)=\begin{pmatrix}\cos\frac\theta2&-\sin\frac\theta2\\\sin\frac\theta2&\cos\frac\theta2\end{pmatrix}$
and $0\le\theta<2\pi$. Alice holds the first qubit ($A$) and Bob the second one ($B$).
\begin{center}
\begin{quantikz}[row sep=0.4cm, column sep=0.4cm]
\lstick{$\ket{q_0}$} & \gate{H} & \ctrl{1} & & \\
\lstick{$\ket{q_1}$} & & \targ{} & \gate{R_y(\theta)} &
\end{quantikz}
\end{center}

\begin{enumerate}
\item Write the matrix of the circuit in the computational basis $\{\ket{00},\ket{01},\ket{10},\ket{11}\}$ \pts{0.5 p}.
\item Write the output state $\ket\psi$ for the input $\ket{00}$. For which values of $\theta$ is it
one of the four Bell states? \pts{0.5 p}
\item Calculate $\rho_A$, $\rho_B$ and the entanglement entropy as a function of $\theta$. Explain
the result \pts{0.5 p}.
\item For $\theta=\pi/2$, find a Schmidt decomposition of $\ket\psi$. Is it unique? Give a second,
different Schmidt decomposition \pts{0.5 p}.
\item We measure the joint observable $Z\otimes Z$ on $\ket\psi$ (general $\theta$). Which are the
possible outcomes, their probabilities and the state after the measurement? Calculate
$\langle Z\otimes Z\rangle$. Is the state still entangled after the measurement? Compare with a
measurement of $Z$ on qubit $B$ only \pts{0.5 p}.
\item The outcome of the $Z\otimes Z$ measurement is not recorded. Write the resulting state
$\rho'$, its purity and its fidelity with $\ket\psi$. Can Alice notice, by measuring her qubit only,
whether the measurement took place? \pts{0.5 p}
\end{enumerate}

\textit{Useful formulas:} $\cos^2\frac x2-\sin^2\frac x2=\cos x$, $\;2\sin\frac x2\cos\frac x2=\sin x$.

%=====================================================================
\ifdefined\nosolutions\else
\newpage
\begin{center}
\Large\textbf{Solutions}
\end{center}
\sol

\section*{1\quad Open question}

The QFT maps $\ket x\mapsto\frac1{\sqrt N}\sum_ye^{2\pi ixy/N}\ket y$, so it acts on a state
$\sum_xa_x\ket x$ by Fourier transforming its \emph{amplitudes}. Its circuit ($H$ gates, controlled
phases $R_d$ and final SWAPs) has $n(n+1)/2=\mathcal O(n^2)$ gates, exponentially fewer than the
$\mathcal O(n2^n)$ operations of the FFT.

This is \textbf{not} an exponential speed-up for classical data, for two reasons:
\begin{itemize}
\item \emph{Input}: the $N$ data must first be loaded as amplitudes. Preparing an arbitrary state
generally needs $\mathcal O(N)$ operations, which already cancels the advantage.
\item \emph{Output}: the Fourier coefficients are amplitudes and cannot be read out. A measurement
gives one value $y$ with probability $|\tilde a_y|^2$; estimating all $N$ coefficients would need
exponentially many repetitions.
\end{itemize}
The QFT is useful as a subroutine when we only need a \emph{global} property encoded in the phases.
In \textbf{phase estimation}, the controlled-$U^{2^j}$ gates (phase kickback) prepare
$\frac1{\sqrt N}\sum_ye^{2\pi i\varphi y}\ket y=\text{QFT}\ket{\varphi N}$. The inverse QFT maps this
state to $\ket{\varphi N}$, so measuring the register gives the eigenphase $\varphi$ with $n$ bits
of precision (exactly if $\varphi N$ is an integer, with probability $\ge4/\pi^2$ otherwise).
(Also valid: \textbf{order finding}, where the QFT turns the periodic state
$\sum_j\ket{x_0+jr}$ into peaks at $y\approx kN/r$, from which $r$ follows by continued fractions.)

\section*{2\quad Short questions}

\begin{enumerate}
\item Yes. The circuit is a multi-controlled $Z$ with controls $q_0$, $q_1$ and target $q_3$, with
$X$ gates before and after on $q_1$. It adds a sign $-1$ only when $q_0=1$, $q_1=0$ and $q_3=1$.
Qubit $q_2$ is not touched, so it can take any value. The gate is diagonal with entries $\pm1$, so it is
a valid phase oracle $U_w\ket x=(-1)^{f(x)}\ket x$, with
\[
\text{marked items: }\ket{w}\in\{\ket{1001},\ket{1011}\}\qquad(M=2,\ N=2^4=16).
\]
Then $\sin\theta_0=\sqrt{M/N}=\sqrt{1/8}=\frac1{2\sqrt2}$, so $\theta_0\approx0.361$ rad $\approx20.7^\circ$.
\[
k_{\max}\simeq\left\lfloor\frac\pi4\sqrt{\frac NM}\right\rfloor=\left\lfloor\frac\pi4\sqrt8\right\rfloor=\lfloor2.22\rfloor=2,
\qquad P_2=\sin^2(5\theta_0)=\sin^2(103.5^\circ)=\frac{121}{128}\approx0.945 .
\]
For comparison, $P_1=\sin^2(3\theta_0)=\frac{25}{32}\approx0.781$. With one extra iteration,
$P_3=\sin^2(7\theta_0)=\sin^2(144.9^\circ)\approx0.33$: the state rotates past $\ket w$ and the
probability drops.

\item
\begin{enumerate}[label=(\roman*)]
\item $U\ket1=e^{3i\pi/4}\ket1=e^{2\pi i\frac38}\ket1$, so $\varphi=\frac38=\frac{x}{2^3}$ and
$x=3\equiv\{011\}$. For $U^\dagger$: $U^\dagger\ket1=e^{-3i\pi/4}\ket1=e^{2\pi i(1-\frac38)}\ket1$, so
$\varphi=\frac58$ and $x=5\equiv\{101\}$. In both cases $\varphi N$ is an integer, so the outcome is
obtained with probability 1.
\item $V\ket1=e^{2\pi i\frac13}\ket1$, so $\varphi=\frac13$ and $\varphi N=\frac83\approx2.67$, which is
not an integer. The most likely outcome is the nearest integer, $x=3\equiv\{011\}$ (i.e.\
$\varphi\approx0.375$), with probability $\ge4/\pi^2\approx0.41$. (The exact value is
$P(x)=\frac{\sin^2(\pi N\delta)}{N^2\sin^2(\pi\delta)}$ with $\delta=\varphi-x/N=-\frac1{24}$, which
gives $P(011)\approx0.69$. The next most likely outcome is $\{010\}$, with $P\approx0.18$.)
\item With $t=3$ bits of precision and $\varepsilon=0.1$:
\[
n=t+\left\lceil\log_2\left(2+\frac1{2\varepsilon}\right)\right\rceil=3+\lceil\log_27\rceil=3+3=6\text{ qubits}.
\]
\end{enumerate}

\item $\vec y_0=\{01001101\}$ gives $y_0=64+8+4+1=77$, with $N=2^8=256$. We look for $k/r$ with
$\left|\frac{y_0}N-\frac kr\right|<\frac1{2N}=\frac1{512}\approx0.00195$. Euclid's algorithm:
$256=3\cdot77+25$, $77=3\cdot25+2$, $25=12\cdot2+1$, $2=2\cdot1$, so
\[
\frac{77}{256}=[0;3,3,12,2],\qquad\text{convergents: }\frac01,\ \frac13,\ \frac3{10},\ \frac{37}{123},\ \frac{77}{256}.
\]
$\left|\frac{77}{256}-\frac13\right|\approx0.033>\frac1{512}$, while
$\left|\frac{77}{256}-\frac3{10}\right|=\frac{1}{1280}\approx0.00078<\frac1{512}$. So $r=10$ (with $r<N$).
Check: $2^1,\dots,2^{10}\bmod11=2,4,8,5,10,9,7,3,6,1$. So $2^{10}\equiv1$, no smaller power gives 1,
and the order is $r=10$.

$\vec y_0^{\,\prime}=\{10000000\}$ gives $y_0'=128$ and $\frac{128}{256}=\frac12=[0;2]$. The
continued fractions suggest $r=2$, but $2^2=4\not\equiv1\pmod{11}$, so this is wrong. The reason is
that $y_0'$ corresponds to $k=5$ and $r=10$: $\frac5{10}=\frac12$, and since $\gcd(5,10)=5>1$ the
fraction reduces and we only obtain a divisor of $r$. We can test multiples of 2:
$2^4\equiv5$, $2^6\equiv9$, $2^8\equiv3$, $2^{10}\equiv1\pmod{11}$, which gives $r=10$. Alternatively,
we repeat the experiment and take the lcm of the denominators obtained.

\end{enumerate}

\section*{3\quad Amplitude damping}

\begin{enumerate}
\item Since $M_1^\dagger=\begin{pmatrix}0&0\\\sqrt\gamma&0\end{pmatrix}$:
\[
E_0=M_0^\dagger M_0=\begin{pmatrix}1&0\\0&1-\gamma\end{pmatrix},\qquad
E_1=M_1^\dagger M_1=\begin{pmatrix}0&0\\0&\gamma\end{pmatrix},\qquad E_0+E_1=\id .
\]
So $\{M_m\}$ is a generalized measurement. Both $E_m$ are positive (diagonal, non-negative entries).
They are not projectors for $0<\gamma<1$: $E_1^2=\gamma E_1\neq E_1$. The operators $M_m$ are not
Hermitian either.

\item With $E_1=\gamma\ket1\!\bra1$:
\[
p(1)=\tr[E_1\rho]=\gamma\rho_{11}=\gamma\,\frac{1-r_z}2,\qquad
\rho_1'=\frac{M_1\rho M_1^\dagger}{p(1)}=\frac{\gamma\rho_{11}\ket0\!\bra0}{\gamma\rho_{11}}=\ket0\!\bra0 .
\]
\[
p(0)=1-\gamma\rho_{11},\qquad
\rho_0'=\frac{M_0\rho M_0^\dagger}{p(0)}=\frac1{1-\gamma\rho_{11}}\begin{pmatrix}\rho_{00}&\sqrt{1-\gamma}\,\rho_{01}\\\sqrt{1-\gamma}\,\rho_{10}&(1-\gamma)\rho_{11}\end{pmatrix}.
\]
Interpretation: a photon can only be emitted if the atom is excited, with probability $\gamma$ times
the population of $\ket1$. If it is detected, the atom is certainly in the ground state $\ket0$. If no
photon is detected, we still learn something: the weight of $\ket1$ decreases (the state is pushed
towards $\ket0$) and the coherences are reduced by $\sqrt{1-\gamma}$.

\item $\rho'=M_0\rho M_0^\dagger+M_1\rho M_1^\dagger=\begin{pmatrix}\rho_{00}+\gamma\rho_{11}&\sqrt{1-\gamma}\,\rho_{01}\\\sqrt{1-\gamma}\,\rho_{10}&(1-\gamma)\rho_{11}\end{pmatrix}$. Then
\[
r_x'=2\,\mathrm{Re}\,\rho'_{01}=\sqrt{1-\gamma}\,r_x,\qquad r_y'=-2\,\mathrm{Im}\,\rho'_{01}=\sqrt{1-\gamma}\,r_y,
\]
\[
r_z'=\rho'_{00}-\rho'_{11}=\rho_{00}-\rho_{11}+2\gamma\rho_{11}=r_z+\gamma(1-r_z)=\gamma+(1-\gamma)r_z .
\]

\item A fixed point needs $r_x=r_y=0$ (for $\gamma>0$) and $r_z=\gamma+(1-\gamma)r_z$, i.e.\ $r_z=1$.
The invariant state is the ground state $\ket0$ (the north pole). After $m$ applications,
$r_{x,y}^{(m)}=(1-\gamma)^{m/2}r_{x,y}\to0$ and $1-r_z^{(m)}=(1-\gamma)^m(1-r_z)\to0$, so every state
decays to $\ket0$.

For $\rho=\id/2$ ($\bm r=\bm0$): $\bm r'=(0,0,\gamma)$, so $\rho'=\frac12(\id+\gamma\sigma_z)\neq\id/2$.
The purity increases from $\frac12$ to $\frac{1+\gamma^2}2$. The Pauli channels only shrink the Bloch
vector ($|2p-1|\le1$) and leave $\id/2$ unchanged (they are unital), so they never increase the
purity. The amplitude damping channel is not unital: it drives the system towards a pure state by
releasing energy (and entropy) to the environment.

\item $\ket+$ has $\bm r=(1,0,0)$, so $\bm r'=(\sqrt{1-\gamma},0,\gamma)$ and
$|\bm r'|^2=1-\gamma+\gamma^2$. Then
\[
\gamma(\rho')=\frac{1+|\bm r'|^2}{2}=1-\frac{\gamma(1-\gamma)}2,\qquad
\lambda_\pm=\frac{1\pm\sqrt{1-\gamma+\gamma^2}}{2}.
\]
The output is pure only for $\gamma=0$ (no decay) or $\gamma=1$ (output $\ket0$). For $\gamma=\frac12$:
$|\bm r'|=\frac{\sqrt3}2$, $\gamma(\rho')=\frac78$, $\lambda_\pm=\frac{2\pm\sqrt3}{4}\approx0.933,\ 0.067$, and
\[
S(\rho')=-0.933\log_20.933-0.067\log_20.067\approx0.093+0.261\approx0.35 .
\]

\item Since $\ket+$ is pure, $F=\bra+\rho'\ket+=\frac{1+\bm r\cdot\bm r'}2=\frac{1+\sqrt{1-\gamma}}{2}$. The
trace distance is
\[
d=\tfrac12|\bm r-\bm r'|=\tfrac12\sqrt{\left(1-\sqrt{1-\gamma}\right)^2+\gamma^2}.
\]
For $\gamma=\frac12$: $F=\frac{1+1/\sqrt2}2\approx0.854$ and $d=\frac12\sqrt{0.086+0.25}\approx0.29$.

For $\ket0$: $\bm r'=\bm r$, so $F=1$ and $d=0$ (the ground state is not affected). For $\ket1$:
$\bm r'=(0,0,2\gamma-1)$, so $F=\frac{1-(2\gamma-1)}2=1-\gamma$ and $d=\frac12|{-1}-(2\gamma-1)|=\gamma$.
For $\gamma=\frac12$ this gives $F=0.5$, $d=0.5$: the excited state is the most affected input.
\end{enumerate}

\section*{4\quad A rotated Bell pair}

Notation: $c\equiv\cos\frac\theta2$, $s\equiv\sin\frac\theta2$.

\begin{enumerate}
\item The circuit is $U=(\id\otimes R_y)\,\text{CNOT}\,(H\otimes\id)$. In $2\times2$ blocks,
$\text{CNOT}\,(H\otimes\id)=\frac1{\sqrt2}\begin{pmatrix}\id&\id\\X&-X\end{pmatrix}$, and
$\id\otimes R_y=\begin{pmatrix}R_y&0\\0&R_y\end{pmatrix}$. With $R_yX=\begin{pmatrix}-s&c\\c&s\end{pmatrix}$:
\[
U=\frac1{\sqrt2}\begin{pmatrix}R_y&R_y\\R_yX&-R_yX\end{pmatrix}
=\frac1{\sqrt2}\begin{pmatrix}c&-s&c&-s\\s&c&s&c\\-s&c&s&-c\\c&s&-c&-s\end{pmatrix}.
\]

\item The output is the first column of $U$:
\[
\ket\psi=\frac1{\sqrt2}\left[\ket0R_y\ket0+\ket1R_y\ket1\right]=\frac1{\sqrt2}\left[c\ket{00}+s\ket{01}-s\ket{10}+c\ket{11}\right].
\]
For $\theta=0$: $\ket\psi=\frac{\ket{00}+\ket{11}}{\sqrt2}=\ket{\beta_{00}}$. For $\theta=\pi$ ($c=0$,
$s=1$): $\ket\psi=\frac{\ket{01}-\ket{10}}{\sqrt2}=\ket{\beta_{11}}$ (the singlet).

\item $\alpha=\frac1{\sqrt2}\begin{pmatrix}c&s\\-s&c\end{pmatrix}$, so
\[
\rho_A=\alpha\alpha^\dagger=\frac12\begin{pmatrix}c^2+s^2&-cs+sc\\-sc+cs&s^2+c^2\end{pmatrix}=\frac\id2,\qquad
\rho_B=\alpha^T\alpha^*=\frac\id2 .
\]
For every $\theta$: $\gamma(\rho_A)=\gamma(\rho_B)=\frac12$ and $S_{EE}=\log_22=1$. The state is
maximally entangled for all $\theta$. The reason is that $\ket\psi=(\id\otimes R_y(\theta))\ket{\beta_{00}}$:
a local unitary acting on a Bell state cannot change its entanglement.

\item For $\theta=\pi/2$, $c=s=\frac1{\sqrt2}$ and $\ket\psi=\frac12\left[\ket{00}+\ket{01}-\ket{10}+\ket{11}\right]$.
Since $\rho_A=\id/2$ is degenerate ($\eta_1^2=\eta_2^2=\frac12$), \emph{any} orthonormal basis of $A$ is
an eigenbasis. Taking $\{\ket0,\ket1\}$:
\[
\bra0_A\ket\psi=\tfrac12(\ket0+\ket1)=\tfrac1{\sqrt2}\ket+,\qquad
\bra1_A\ket\psi=\tfrac12(-\ket0+\ket1)=\tfrac1{\sqrt2}\left(-\ket-\right),
\]
\[
\boxed{\ket\psi=\frac1{\sqrt2}\left[\ket0\ket+-\ket1\ket-\right]}
\]
The decomposition is \textbf{not unique}, because the Schmidt coefficients are degenerate. Taking
$\{\ket+,\ket-\}$ for $A$ instead:
\begin{align*}
\bra+_A\ket\psi&=\tfrac1{2\sqrt2}\left[(\ket0+\ket1)+(-\ket0+\ket1)\right]=\tfrac1{\sqrt2}\ket1,\\
\bra-_A\ket\psi&=\tfrac1{2\sqrt2}\left[(\ket0+\ket1)-(-\ket0+\ket1)\right]=\tfrac1{\sqrt2}\ket0,
\end{align*}
so $\ket\psi=\frac1{\sqrt2}\left[\ket+\ket1+\ket-\ket0\right]$ is another valid Schmidt decomposition.
In both cases $\eta_1=\eta_2=\frac1{\sqrt2}$ and $S_{EE}=1$.

\item $Z\otimes Z$ has eigenvalues $+1$ (even parity) and $-1$ (odd parity), with projectors
$P_+=\ket{00}\!\bra{00}+\ket{11}\!\bra{11}$ and $P_-=\ket{01}\!\bra{01}+\ket{10}\!\bra{10}$.
\begin{align*}
P_+\ket\psi&=\tfrac c{\sqrt2}(\ket{00}+\ket{11}), & p(+1)&=c^2=\cos^2\tfrac\theta2, & \ket{\psi_+}&=\ket{\beta_{00}};\\
P_-\ket\psi&=\tfrac s{\sqrt2}(\ket{01}-\ket{10}), & p(-1)&=s^2=\sin^2\tfrac\theta2, & \ket{\psi_-}&=\ket{\beta_{11}}.
\end{align*}
$\langle Z\otimes Z\rangle=c^2-s^2=\cos\theta$. Both post-measurement states are Bell states, so the
state is still maximally entangled. A joint (non-local) observable can preserve entanglement. On the
other hand, measuring $Z$ on $B$ only, with $P_m=\id\otimes\ket m\!\bra m$, leaves a product state
$\ket{\psi'_A}\ket m$ and destroys the entanglement. (For example, outcome $0$:
$p=\frac12$, $\ket{\psi'_A}=c\ket0-s\ket1$.)

\item If the outcome is not recorded:
\[
\rho'=\sum_\pm P_\pm\ket\psi\!\bra\psi P_\pm=c^2\ket{\beta_{00}}\!\bra{\beta_{00}}+s^2\ket{\beta_{11}}\!\bra{\beta_{11}}.
\]
The two Bell states are orthogonal, so the eigenvalues of $\rho'$ are $c^2,s^2$ and
\[
\gamma(\rho')=c^4+s^4=1-\tfrac12\sin^2\theta .
\]
It is pure only for $\theta=0,\pi$, when $\ket\psi$ is already an eigenstate of $Z\otimes Z$. Since
$\ket\psi$ is pure, $F=\bra\psi\rho'\ket\psi=c^2|\braket{\beta_{00}|\psi}|^2+s^2|\braket{\beta_{11}|\psi}|^2$,
with $\braket{\beta_{00}|\psi}=c$ and $\braket{\beta_{11}|\psi}=s$:
\[
F(\ket\psi\!\bra\psi,\rho')=c^4+s^4=1-\tfrac12\sin^2\theta,
\]
which is minimal ($F=\frac12$) for $\theta=\pi/2$.
Alice cannot notice anything: $\tr_B\ket{\beta_{00}}\!\bra{\beta_{00}}=\tr_B\ket{\beta_{11}}\!\bra{\beta_{11}}=\id/2$,
so $\rho'_A=c^2\frac\id2+s^2\frac\id2=\frac\id2=\rho_A$. Her local statistics are exactly the same as
before the measurement.
\end{enumerate}
\fi

\end{document}
